\documentclass{article}

\usepackage{nips15submit_e,times}
\usepackage{hyperref}
\usepackage{xurl}
\usepackage{url}
\usepackage{amsmath}
\usepackage{amssymb}
\usepackage{graphicx}
\usepackage{booktabs}
\usepackage{multirow}
\usepackage{float}
\usepackage{natbib}
\usepackage{microtype}
\usepackage{subcaption}
\hypersetup{hidelinks}
\raggedbottom
\setlength{\textfloatsep}{10pt plus 2pt minus 2pt}
\setlength{\floatsep}{8pt plus 2pt minus 2pt}
\setlength{\intextsep}{8pt plus 2pt minus 2pt}
\setlength{\abovecaptionskip}{4pt}

\newcommand{\blankcell}{\textemdash}
\input{../results/table_rows.tex}

\title{A Concise FlashAttention-1 Forward Kernel in CUDA C++}

\author{
Calder Katyal \\
Yale University \\
CPSC 4240: Parallel Programming \\
\texttt{calder.katyal@yale.edu}
\And
Devon Sawyer \\
Yale University \\
CPSC 4240: Parallel Programming \\
\texttt{devon.sawyer@yale.edu}
\And
Akshat Badgamia \\
Yale University \\
CPSC 4240: Parallel Programming \\
\texttt{akshat.badgamia@yale.edu}
}

\nipsfinalcopy

\begin{document}

\maketitle

\begin{abstract}
FlashAttention accelerates the most important operation in modern machine learning--attention--by reducing traffic between high-bandwidth memory (HBM) and on-chip SRAM. This final project provides a small, compact, CUDA-C++ version of the FlashAttention-1 forward pass, which retains the major algorithmic concepts identified in the FlashAttention-1 paper and source code: exact online softmax, tiled on-chip working sets, fused value accumulation, mixed precision, and tensor core-based scoring. The CUDA-C++ version has been constructed with the objective of removing most of the extra features and other low-level pipelining mechanisms that are found in the production package and that make it harder to read. On the benchmarks that were run during this work, the resulting kernel is approximately $17.49\times$ faster than the PyTorch baseline when running sequences of length $4096$ ($25.35\times$ for official FlashAttention-1 at the same sequence length). When ablating individual components within this implementation to identify their contributions to the improved performance, our results showed that tensor cores contribute the most to performance, with vectorized tile movement, online softmax, and shared-memory tiling all contributing moderately. Additionally, because the full score matrix is never materialized, memory usage is also dramatically lower than in the PyTorch baseline.\footnote{View our codebase at \url{https://github.com/calderkatyal/CPSC_4240_Final_Project}.}
\end{abstract}

\section{Introduction}
Self-attention is a central primitive in transformer models, but the naive implementation is expensive in both computation and memory traffic. Given queries $Q \in \mathbb{R}^{N \times d}$, keys $K \in \mathbb{R}^{N \times d}$, and values $V \in \mathbb{R}^{N \times d}$, standard attention computes
\begin{equation}
\operatorname{Attn}(Q,K,V) = \operatorname{softmax}\left(\frac{QK^\top}{\sqrt{d}}\right)V.
\end{equation}
Attention (and in particular self-attention) is one of the most important primitives within transformer models. However, self-attention has two primary costs in terms of memory access patterns and computational overhead. The first is the $O(N^2)$ computational cost associated with calculating the scores of each pair of tokens using $QK^\top$. The second is the memory access cost of materializing the full $N\times N$ attention matrix in high-bandwidth memory (HBM). This problem is addressed through FlashAttention as described in \citet{dao2022flashattention}. FlashAttention rearranges how exact attention operates, creating three distinct operations: score computation, normalization, and value accumulation. Each of these is tiled over a subset of the total number of sequences being processed, resulting in no need to explicitly write out the entire $N\times N$ attention matrix.

Therefore, instead of implementing the complete \texttt{flash\_attn} package, our main focus is on developing an efficient forward pass based upon the FlashAttention-1 paper. The central questions are:
\begin{enumerate}
    \item How much of FlashAttention-1's forward-pass advantage can be recovered with a concise implementation that keeps only the main ideas?
    \item Which of those ideas matter most in practice on the benchmarked workload?
\end{enumerate}
To answer these questions, the benchmark suite compares our FA, four pure ablations, a shared PyTorch baseline, and the official FlashAttention-1 forward kernel on the same machine.

\section{Background}
\subsection{Memory traffic}
For a single head, standard attention contains three conceptual stages: score formation, row-wise normalization, and value mixing. In matrix form,
\begin{align}
S &= \frac{QK^\top}{\sqrt{d}}, \\
P &= \operatorname{softmax}(S), \\
O &= PV.
\end{align}
The basic formulation for attention computation is simple mathematically; however, in terms of actual implementation on a GPU, it has an extremely high cost due to reading and writing the entire $N\times N$ matrix many times. Specifically, this cost becomes extreme when computing sequences with a large number of positions, since the size of the matrix being read from and written to will dominate memory access and intermediate results. FlashAttention's key contribution was to reformulate exact attention so that, rather than maintaining the full matrix of attention scores across all rows, it computes each row separately while keeping only enough per-row normalization information to preserve exactness \citep{dao2022flashattention}.

\subsection{Online softmax}
A central algebraic trick in FlashAttention-1 is to update the softmax normalization across the rows by means of an incremental method that does not involve approximations. Assume that at present a given row has already computed over a subset of the keys it will need to process, and therefore has a current running maximum value of $m$, along with a running denominator value of $\ell = \sum_{j} e^{s_j-m}$. At this point, we have a new set of values for the row; they are denoted by the symbol $\hat{m}$ and represent the largest element within the newly received tile of scores. We now take the new row maximum value as $m'=\max(m,\hat{m})$. Therefore we can express the new row denominator precisely using
\begin{equation}
\ell' = e^{m - m'} \ell + \sum_{j \in \text{new tile}} e^{s_j - m'}.
\end{equation}
The same rescaling factor $e^{m - m'}$ may also be applied to the running sum (the accumulator) before adding the value of the current tile. It is this iterative process that allows FlashAttention to continue to operate exactly while discarding the old score tile once its contribution has been incorporated into the overall state of the computation \citep{dao2022flashattention}.

Therefore, this method of decreasing memory usage does not modify the actual attention operation; instead, it simply changes the way that the computation is ordered in terms of time.

\subsection{Later variations}

FlashAttention-2 incorporates improvements over the original version of FlashAttention regarding forward pass workload partitioning, and also has improved performance relative to non-matmul operations \citep{dao2023flashattention2}. FA2 provides a strong example of how improving both partitioning of workload among threads/warps and reducing the amount of shared processing needed among them can result in greater speed.

FlashAttention-3 represents an additional architectural improvement targeted at Hopper class systems \citep{shah2024flashattention3}. Its primary advancements include: (i) asynchronous execution; (ii) warp specialization; and (iii) increased utilization of lower precision hardware capabilities. Therefore, this variation represents a more detailed level of integration between the FlashAttention kernel and system hardware than either of the first two.

FlashAttention-4 continues this process of detailed kernel/hardware integration for newer systems \citep{zadouri2026flashattention4}. Its focus includes: (i) large-scale rework of pipelines; (ii) greater degrees of overlap during matmul execution; and (iii) elimination of remaining bottlenecks associated with non-matmul operations such as exponential calculations. As such, it represents an even greater departure from the scope of a single CUDA kernel than any of the previous versions.

As our implementation is forward-only, we prioritize the parts of FA1 that most directly shape forward-pass performance:
\begin{itemize}
    \item no explicit score-matrix materialization,
    \item shared-memory tiling,
    \item online softmax,
    \item fused exact attention,
    \item mixed precision, and
    \item tensor-core-friendly score computation.
\end{itemize}
These are the parts of FlashAttention that are both conceptually central and broadly relevant on modern NVIDIA GPUs.

\begin{figure}[H]
    \centering
    \fbox{
        \parbox{0.94\linewidth}{
            \centering
            \small
            \textbf{The PyTorch baseline materializes the score matrix, whereas FlashAttention keeps only a working set}\\[5pt]
            \begin{minipage}[t]{0.46\linewidth}
                \centering
            \textbf{PyTorch baseline}\\[3pt]
                \fbox{\parbox{0.8\linewidth}{\centering Read $Q,K$}}\\[3pt]
                $\Downarrow$\\[3pt]
                \fbox{\parbox{0.8\linewidth}{\centering Form full $N\times N$ score matrix $S$}}\\[3pt]
                $\Downarrow$\\[3pt]
                \fbox{\parbox{0.8\linewidth}{\centering Write $S$ to HBM}}\\[3pt]
                $\Downarrow$\\[3pt]
                \fbox{\parbox{0.8\linewidth}{\centering Read $S$ back for softmax}}\\[3pt]
                $\Downarrow$\\[3pt]
                \fbox{\parbox{0.8\linewidth}{\centering Read $V$, produce $O$}}
            \end{minipage}\hfill
            \begin{minipage}[t]{0.46\linewidth}
                \centering
                \textbf{FlashAttention-style execution}\\[3pt]
                \fbox{\parbox{0.8\linewidth}{\centering Stage one query block and one KV block in SRAM}}\\[3pt]
                $\Downarrow$\\[3pt]
                \fbox{\parbox{0.8\linewidth}{\centering Compute score subtiles}}\\[3pt]
                $\Downarrow$\\[3pt]
                \fbox{\parbox{0.8\linewidth}{\centering Update online softmax statistics}}\\[3pt]
                $\Downarrow$\\[3pt]
                \fbox{\parbox{0.8\linewidth}{\centering Accumulate the output block immediately}}\\[3pt]
                $\Downarrow$\\[3pt]
                \fbox{\parbox{0.8\linewidth}{\centering Write only final $O$}}
            \end{minipage}
        }
    }
    \caption{The core memory-efficiency idea of FlashAttention. The PyTorch baseline repeatedly touches the full score matrix in high-bandwidth memory, whereas FlashAttention keeps only the currently active query/KV block and per-row running statistics live at any moment.}
    \label{fig:memory_view}
\end{figure}

\section{Methods}
\subsection{PyTorch baseline}
The reference baseline is equivalent to the PyTorch attention path described in the official FlashAttention benchmark script. It creates the complete score matrix using a batched matmul, adds causal masks, applies row-wise softmax, and finally multiplies the resulting attention probabilities by $V$. The input data type of the baseline is FP16 to match our implementation as well as the official FA1. Note that the PyTorch baseline creates the complete score matrix, making it useful for benchmarking.

\subsection{Our FA}
Our FA mirrors the structure of the forward algorithm in the FA1 paper and repository \citep{dao2022flashattention,dao2022flashattentionrepo}. Each block of threads has ownership of a specific query block and moves through the sequence in K/V tiles. During benchmarking we configure each block to process $64$ query rows at a time and move through $128$-row K/V tiles. All queries and the currently processed K/V tile are placed on-chip, all subtiles of scores are created using WMMA tensor core fragments, and after every K/V tile the running softmax state is updated. We use $m_i$ and $\ell_i$ to represent the running row maxima and running normalizing terms for query row $i$. For a new K/V tile producing scores $s_{ij}$, the online recurrence is
\begin{align}
m_i' &= \max\left(m_i, \max_j s_{ij}\right), \\
\ell_i' &= e^{m_i - m_i'} \ell_i + \sum_j e^{s_{ij} - m_i'}, \\
o_i' &= e^{m_i - m_i'} o_i + \sum_j e^{s_{ij} - m_i'} V_j .
\end{align}
After the final KV block, each output row is normalized by $\ell_i'$. This is the main algorithmic reason FlashAttention can remain exact while avoiding the full $N\times N$ attention matrix \citep{dao2022flashattention}.

The implementation is intentionally smaller than the official package, but our FA preserves the core forward-pass structure:
\begin{itemize}
    \item a query block and the current key/value block are staged in shared memory;
    \item the score computation is built from WMMA tensor-core subtiles inside the staged K tile;
    \item softmax normalization and value accumulation are fused into the same kernel; and
    \item only per-row running state and the final output are kept live between KV blocks, without materializing the full score matrix.
\end{itemize}

Notably, our FA does not include dropout, backward kernels, variable-length batching, MQA/GQA, KV cache support, BF16 variants, or the large family of architecture-specific kernels used in the package. We also omitted some of the more advanced scheduling mechanisms from the source tree, namely more complex forms of software pipelining and more specialized uses of shared memory. The omission of those elements keeps the code much smaller while still preserving a similar forward attention rule.

\begin{figure}[H]
    \centering
    \fbox{
        \parbox{0.94\linewidth}{
            \centering
            \small
            \textbf{HBM stores the full tensors, while SRAM stores only the current working set}\\[5pt]
            \textbf{HBM (global memory)}\\[3pt]
            \fbox{\parbox{0.18\linewidth}{\centering full $Q$\\tensor}} \hfill
            \fbox{\parbox{0.18\linewidth}{\centering full $K$\\tensor}} \hfill
            \fbox{\parbox{0.18\linewidth}{\centering full $V$\\tensor}} \hfill
            \fbox{\parbox{0.18\linewidth}{\centering final $O$\\tensor}}\\[5pt]
            $\Downarrow$ \hspace{3pt} \textit{one thread block loads only the active tiles} \hspace{3pt} $\Downarrow$\\[5pt]
            \textbf{SRAM / on-chip working set for one block}\\[3pt]
            \fbox{\parbox{0.18\linewidth}{\centering $Q_i$ tile\\$64 \times d$}} \hfill
            \fbox{\parbox{0.18\linewidth}{\centering $K_j$ tile\\$128 \times d$}} \hfill
            \fbox{\parbox{0.18\linewidth}{\centering $V_j$ tile\\$128 \times d$}} \hfill
            \fbox{\parbox{0.18\linewidth}{\centering running $(m,\ell,o)$\\per active row}}\\[5pt]
            \fbox{\parbox{0.82\linewidth}{\centering Ephemeral score and probability subtiles are computed, consumed immediately by the online update, and then discarded.}}\\[4pt]
            \fbox{\parbox{0.82\linewidth}{\centering No full $N\times N$ score matrix is stored in either HBM or SRAM.}}
        }
    }
    \caption{Explicit memory-hierarchy view of the FA1 kernel used in the benchmark. The full tensors live in HBM, but the kernel keeps only the current query tile, the current K/V tile, and the row-wise running state on chip.}
    \label{fig:hbm_sram_map}
\end{figure}

\subsection{Ablations}
We include four ablation kernels so the results can isolate which FA1 ingredients matter most.
\begin{itemize}
    \item \textbf{No tensor cores.} This kernel retains the same staged query/KV tile layout, online-softmax recurrence, and fused forward pass as the original; however, it uses scalar per-thread math to replace the WMMA score and update paths.
    \item \textbf{No vectorized loads.} This kernel retains the same core algorithm; however, it disables the 16-byte Q/K/V tile loads that are used in the primary path.
    \item \textbf{No online softmax.} This kernel retains the same outer-block query/KV structure, vectorized tile movement, tensor-core score path, and WMMA value-update path as the original; however, it uses a two-pass computation based on previously computed row maxima to replace the online recurrence.
    \item \textbf{No SRAM tiling.} This kernel retains the same outer-block structure, tensor-core score path, WMMA value-update path, and exact online-softmax logic as the original; however, it loads K/V MMA micro-tiles directly from global memory rather than staging them in shared-memory K/V blocks.
\end{itemize}

\section{Experimental Setup}
We study dense causal self-attention in a controlled forward-only setting. All methods use the same tensor shapes and masking convention so that runtime differences reflect kernel organization rather than changes in the mathematical task. Our FA and the official baseline both run in mixed precision with FP16 inputs. Output accumulation and the reference comparison are performed in FP32, which matches the numerically stable regime used in the FlashAttention literature \citep{dao2022flashattention}.

\begin{table}[H]
    \centering
    \small
    \caption{Benchmark configuration used by the CUDA benchmark suite.}
    \label{tab:benchcfg}
    \begin{tabular}{lr|lr}
        \toprule
        Quantity & Value & Quantity & Value \\
        \midrule
        Batch size $B$ & 2 & Heads $H$ & 8 \\
        Head dimension $d$ & 64 & Causal mask & yes \\
        Input dtype & FP16 & Accumulation dtype & FP32 \\
        Sequence lengths & 128--4096 & Warmup iters & 10 \\
        Timed iters & 50 & Official package dtype & FP16 \\
        Head-dim support & $\{32, 64, 128\}$ & Target hardware & NVIDIA RTX 5000 Ada Generation  \\
        \bottomrule
    \end{tabular}
\end{table}

The benchmark fixes batch size, head count, head dimension, and masking pattern, and varies only sequence length. This makes the role of memory traffic especially visible, since increasing $N$ enlarges the fully materialized score matrix quadratically while leaving the head dimension fixed. The main sweep uses $B=2$, $H=8$, $d=64$, causal masking, and sequence lengths from $128$ to $4096$. We report runtime, peak GPU memory, and maximum absolute error relative to an FP32 reference.

The evaluation reports three classes of comparisons.
\begin{enumerate}
    \item \textbf{PyTorch baseline vs.\ our FA.} This isolates the main FlashAttention-1 design choices inside one compact CUDA kernel.
    \item \textbf{Our FA vs.\ official FlashAttention-1.} This measures the gap between the compact implementation and the production FA1 kernel.
    \item \textbf{FA1 ablations.} These isolate which of the core FA1 ingredients are most important in practice.
\end{enumerate}

The official baseline anchors the comparison to the original IO-aware reformulation and provides a reference point for how much performance remains tied to the heavier engineering in the package kernel.

\section{Results}

\begin{figure}[H]
    \centering
    \begin{subfigure}[t]{0.48\columnwidth}
        \centering
        \includegraphics[width=\linewidth]{../results/gpu_comparison_runtime.pdf}
    \end{subfigure}\hfill
    \begin{subfigure}[t]{0.48\columnwidth}
        \centering
        \includegraphics[width=\linewidth]{../results/gpu_comparison_speedup.pdf}
    \end{subfigure}

    \vspace{0.12in}

    \begin{subfigure}[t]{0.48\columnwidth}
        \centering
        \includegraphics[width=\linewidth]{../results/gpu_comparison_memory.pdf}
    \end{subfigure}\hfill
    \begin{subfigure}[t]{0.48\columnwidth}
        \centering
        \includegraphics[width=\linewidth]{../results/gpu_comparison_ablation.pdf}
    \end{subfigure}
    \caption{Runtime, speedup, memory, and ablation views for the merged benchmark suite.}
    \label{fig:gpuplots}
\end{figure}

\begin{table}[H]
    \centering
    \small
    \caption{Main runtime comparison. Entries report average kernel runtime in milliseconds.}
    \label{tab:runtime}
    \begin{tabular}{lcccc}
        \toprule
        Method & $N=512$ & $N=1024$ & $N=2048$ & $N=4096$ \\
        \midrule
        \runtimeTableRows
        \bottomrule
    \end{tabular}
\end{table}

\begin{table}[H]
    \centering
    \small
    \caption{Speedup relative to the shared PyTorch attention baseline. This table directly quantifies the effect of the FlashAttention-1 core under one common denominator for all methods.}
    \label{tab:speedups}
    \begin{tabular}{lcccc}
        \toprule
        Method & $N=512$ & $N=1024$ & $N=2048$ & $N=4096$ \\
        \midrule
        \speedupTableRows
        \bottomrule
    \end{tabular}
\end{table}

\begin{table}[H]
    \centering
    \small
    \caption{Ablation table for the FlashAttention-1 core at $N=2048$. The speedup column is relative to the shared PyTorch attention baseline at the same sequence length.}
    \label{tab:ablation}
    \begin{tabular}{lccc}
        \toprule
        Method & Runtime at $N=2048$ (ms) & Speedup vs.\ PyTorch & Max abs.\ error \\
        \midrule
        \ablationTableRows
        \bottomrule
    \end{tabular}
\end{table}

Figure~\ref{fig:gpuplots} and Tables~\ref{tab:runtime}--\ref{tab:ablation} together show that the compact kernel captures much of FA1's speed while remaining short of the full production implementation.

\subsection{Runtime and speedup}
The main comparison is the PyTorch baseline versus our FA. At sequence length $N=512$, our FA was $4.07\times$ faster than the PyTorch baseline. By $N=1024$, it was $5.70\times$ faster, by $N=2048$, $12.15\times$ faster, and by $N=4096$, $17.49\times$ faster. This is what we predicted based on the FA1 argument in the paper: as the sequence length increases, the PyTorch baseline gets progressively worse at materializing and then reading the entire score matrix, while the working set of our FA remains bounded by the query tile, the current K/V tile, and the row-wise running state.

Memory usage follows a similar pattern. At $N=4096$, the PyTorch baseline peaks at about $1.28$~GB, while our FA uses about $33.6$~MB. That is a reduction of roughly $38\times$. The difference is much smaller at shorter sequence lengths, but as $N$ increases, so does the size of the PyTorch baseline score matrix and therefore the gap in memory usage.

\subsection{Comparison with official FA1}
As stated earlier, our FA is very competitive with the official baseline for both short and moderate-length sequences. At sequence length $N=512$, our FA is actually slightly faster than official FA1. That should not be read as a claim that our FA is generally better. The benchmark fixes one shape and one narrow feature set, whereas the package kernel is designed to cover a broader range of head dimensions, masks, and deployment settings.

At longer sequence lengths, the production kernel pulls ahead in the expected way. At $N=2048$, our FA runs in $0.1373$~ms, compared with $0.0968$~ms for official FA1. At $N=4096$, the gap widens to $0.4435$~ms versus $0.3060$~ms for official FA1. In speedup terms, our FA reaches $17.49\times$ versus PyTorch at $N=4096$, while official FA1 reaches $25.35\times$. It appears most likely that the compact kernel has captured the major algorithmic advantages of FA1; however, it still lacks some of the heavier engineering that matters when the sequence becomes sufficiently long that the package kernel can amortize its additional machinery.

\subsection{Ablation results}
The ablations confirm that each feature contributes to the full gain, although one factor is significantly more important. At $N=2048$, removing tensor cores increases runtime from $0.1373$~ms to $1.2214$~ms, making this the largest single drop. This is consistent with the fact that the score calculation is still the dominant arithmetic kernel even after the IO reorganization.

The results from the other ablations provide additional information as well. Ablating the use of vectorized Q/K/V loads increased runtime to $0.3990$~ms. This demonstrates that the memory access path is important, even if one performs tiling at a high level. Also, replacing the online softmax operation with the two-pass exact version increased the runtime to $0.2841$~ms. This indicates that the online recurrence is not merely a theoretical mechanism to achieve exactness, but that there is also a significant practical benefit to reducing the number of operations per row in this case. Finally, when disabling the K/V shared-memory tiling we see runtime increase to $0.2405$~ms. While still much faster than PyTorch, since the kernel performs exact fused tile-wise computation and never computes the entire score matrix, losing this tiling is still quite costly compared with running the full FA1 kernel.

\section{Discussion}
The results support a simple conclusion: the primary concepts for a forward pass in the paper are sufficient to realize a significant portion of FA1's practical advantages. The forward pass implemented in this work retains the essential algorithmic components of FA1--tile-by-tile exact attention, online softmax, and immediate value accumulation.

At the same time, the comparison with the official FA1 model clearly identifies the point at which brevity ceases to provide additional efficiency. The higher-level optimization and tuning used in the development of the package kernel produce greater benefits in terms of execution time than what was possible using our approach. This difference is especially evident when working with longer input sequences, in which the optimized package consistently outperforms our implementation.

\section{Limitations}
This implementation still has several deliberate limitations.
\begin{itemize}
    \item The code is forward-only and does not include backward or dropout kernels.
    \item The kernels currently specialize the common head dimensions $\{32, 64, 128\}$ rather than the full set of multiples of $16$.
    \item The implementation uses FP16 inputs and FP32 accumulation, but does not yet support BF16.
    \item The comparison against official FA1 is numerically meaningful and practically useful, but it is not a full feature-for-feature comparison.
\end{itemize}
These limitations are acceptable for the stated goal of this report, as our emphasis is on understanding the FA1 forward algorithm and measuring the impact of its main ingredients--not on reproducing the entire production stack.

\section{Conclusion}
In this project, we saw that FlashAttention-1 can largely be captured in a minimalistic, yet efficient, CUDA kernel implementation without sacrificing its primary performance characteristics. We find that by implementing a much reduced version of the FlashAttention-1 concepts, significant speed improvements over the PyTorch baseline were achieved. Additionally, our version significantly decreased memory usage and remained only moderately farther away from the optimized FA1 package kernel at longer input sequences.
Further ablation experiments demonstrated that tensor core operations contributed the greatest amount to the overall speed improvement; however, vectorized movement, online softmax, and shared-memory tiling also made independent, though smaller, contributions to achieving this result.
Together, these results validate the assertion that the fundamental ideas behind FA1 are both understandable and effective in a much more limited body of source code.

\clearpage

\begin{thebibliography}{4}

\bibitem[Dao et~al.(2022)Dao, Fu, Ermon, Rudra, and
R{\'e}]{dao2022flashattention}
Tri Dao, Daniel Y. Fu, Stefano Ermon, Atri Rudra, and Christopher R{\'e}.
\newblock FlashAttention: Fast and memory-efficient exact attention with IO-awareness.
\newblock In \emph{Advances in Neural Information Processing Systems}, 2022.

\bibitem[Dao(2022)]{dao2022flashattentionrepo}
Tri Dao.
\newblock FlashAttention \texttt{v1.0.9} source tree.
\newblock GitHub repository, 2022.

\bibitem[Dao(2023)]{dao2023flashattention2}
Tri Dao.
\newblock FlashAttention-2: Faster attention with better parallelism and work partitioning.
\newblock 2023.
\newblock arXiv:2307.08691.

\bibitem[Shah et~al.(2024)Shah, Fu, Kumbong, Roumeliotis, and
Dao]{shah2024flashattention3}
Jay Shah, Ganqu Cui Fu, Hazael Q. Kumbong, Konstantinos Roumeliotis, and Tri Dao.
\newblock FlashAttention-3: Fast and accurate attention with asynchronous low-precision computation.
\newblock In \emph{Advances in Neural Information Processing Systems}, 2024.

\bibitem[Zadouri et~al.(2026)Zadouri, Hoehnerbach, Shah, Liu, Thakkar, and
Dao]{zadouri2026flashattention4}
Ted Zadouri, Markus Hoehnerbach, Jay Shah, Timmy Liu, Vijay Thakkar, and Tri Dao.
\newblock FlashAttention-4: Algorithm and kernel pipelining co-design for asymmetric hardware scaling.
\newblock 2026.
\newblock arXiv:2603.05451.

\end{thebibliography}

\end{document}
